%==============================================================================
% CORPUS CONSTRUCTION tables (§3.3). Paste where referenced.
% booktabs is loaded in the AAAI template.
%
% COUNTS FINAL: recounted from results/final_status.jsonl (10,474 records, 2026-07-14).
% Status map -> AUSTIN_PROVEN 262, NO_FINITE_MODEL 4141, TRIVIAL 3080,
% HAS_FINITE_MODEL 1042, SATISFIABLE_ONLY 43, OPEN 1906.
%==============================================================================


%------------------------------------------------------------------------------
% Table C1 — corpus composition. Single column; \begin{table}.
% Grouped by what the pipeline does with each candidate.
%------------------------------------------------------------------------------
\begin{table}[t]
\centering
\small
\begin{tabular}{lr}
\toprule
Candidate outcome & Count \\
\midrule
\multicolumn{2}{l}{\emph{Admitted (certified no nontrivial finite model)}} \\
\quad Austin, model exhibited          & 262 \\
\quad Admissible, unresolved (hard tier) & 4{,}141 \\
\quad Trivial ($L \models x{=}y$)      & 3{,}080 \\
\midrule
\multicolumn{2}{l}{\emph{Discarded / uncertified}} \\
\quad Has a nontrivial finite model    & 1{,}042 \\
\quad Nontrivial model, admissibility uncertified & 43 \\
\quad Not certified admissible         & 1{,}906 \\
\midrule
Total candidates                       & 10{,}474 \\
\bottomrule
\end{tabular}
\caption{Composition of the generated pool. A candidate is \emph{admitted} only on a
proof that it has no nontrivial finite model; admitted laws are each trivial or Austin,
and the \emph{hard tier} is the admissible subset the baseline leaves unresolved.
Counts are by distinct law string; modulo equivalence the corpus is smaller
(Table~\ref{tab:classes}).}
\label{tab:corpus}
\end{table}


%------------------------------------------------------------------------------
% Table C2 — distinctness / class collapse. Single column; \begin{table}.
% From equiv_sample on the Austin-proven laws.
%------------------------------------------------------------------------------
\begin{table}[t]
\centering
\small
\begin{tabular}{lr}
\toprule
Distinctness (Austin-proven sample) & \\
\midrule
Laws compared                        & 262 \\
Confirmed-equivalent pairs           & 255 \\
Net class merges                     & 67 \\
\midrule
Distinct classes (upper bound)       & $\leq 195$ \\
Collapse                             & $\approx 26\%$ \\
\bottomrule
\end{tabular}
\caption{Distinctness on the Austin-proven laws. Pairs are \emph{separated} by an
explicit model in which one law holds and the other fails---a certificate that they are
distinct---so the class count is anchored from below; unresolved pairs are counted as
distinct, making $195$ an upper bound. The $262$ laws collapse to $\leq 195$ classes
($67$ net merges from $255$ confirmed-equivalent pairs). As an illustration of how sharply
the separator distinguishes laws, a $48$-law pilot separated $1{,}123$ of its $1{,}128$
pairs. The corpus is reported in classes, not laws.}
\label{tab:classes}
\end{table}


%------------------------------------------------------------------------------
% Table C3 (OPTIONAL) — yield of extension vs random sampling.
% Include only if you have the numbers; it makes the "extension, not sampling"
% claim quantitative. Single column.
%------------------------------------------------------------------------------
\begin{table}[t]
\centering
\small
\begin{tabular}{lcc}
\toprule
Generation method & Admissible rate & Rel. yield \\
\midrule
Random order-$\geq 6$ sampling & $\sim\!0.?\%$    & $1\times$ \\
One-operation extension        & $\sim\!??\%$     & $\sim\!500\times$ \\
\bottomrule
\end{tabular}
\caption{Yield of admissible laws under random sampling versus one-operation extension
from Austin seeds; extension raises the admissible rate by about two orders of
magnitude, which is what makes bulk generation practical.
% FILL both rates from the generation logs; the ~500x figure is from the harvest notes.}
\label{tab:yield}
\end{table}
